% note.tex  --  SLE Hadamard states on Bianchi VII_0 with Bieberbach spatial slices
% Agent A4, 2026-05-03.  Draft for Kévin Remondière; iterate before submission.
% Target journal: Letters in Mathematical Physics / Communications in Mathematical Physics.
%
% VERIFIED CITATIONS (arXiv API, 2026-05-03):
%   [BN23]  Banerjee-Niedermaier, arXiv:2305.11388, JMP 64 (2023) 113503
%   [AV13]  Avetisyan-Verch,      arXiv:1212.6180,  CMP (2013)
%   [TB13]  Them-Brum,            arXiv:1302.3174,  CQG 30 (2013) 235035
%   [Rad96] Radzikowski,          CMP 179 (1996) 529-553
%   [Olb07] Olbermann,            arXiv:0704.2986
%   [BFV03] Brunetti-Fredenhagen-Verch, CMP 237 (2003) 31-68, arXiv:math-ph/0112041
%   [Ver94] Verch, CMP 160 (1994) 507-536
%
% NOT YET FOUND on arXiv (need library check before submission):
%   - "Brum-Them 2013" SLE Hadamard estimates: arXiv:1302.3174 authors are THEM-BRUM
%     (order reversed from task description).  Confirmed: Kolja Them & Marcos Brum.

\documentclass[12pt, a4paper]{article}

\usepackage{amsmath, amsthm, amssymb, mathrsfs}
\usepackage{hyperref}
\usepackage{geometry}
\geometry{margin=2.5cm}

\newtheorem{theorem}{Theorem}[section]
\newtheorem{proposition}[theorem]{Proposition}
\newtheorem{lemma}[theorem]{Lemma}
\newtheorem{corollary}[theorem]{Corollary}
\newtheorem{remark}[theorem]{Remark}
\newtheorem{definition}[theorem]{Definition}

\theoremstyle{remark}
\newtheorem{gap}[theorem]{Gap}

% Notation shortcuts
\newcommand{\bR}{\mathbb{R}}
\newcommand{\bZ}{\mathbb{Z}}
\newcommand{\bC}{\mathbb{C}}
\newcommand{\cH}{\mathcal{H}}
\newcommand{\cA}{\mathcal{A}}
\newcommand{\cF}{\mathcal{F}}
\newcommand{\cS}{\mathcal{S}}
\newcommand{\SLE}{\mathrm{SLE}}
\newcommand{\WF}{\mathrm{WF}}
\newcommand{\supp}{\mathrm{supp}}
\newcommand{\tr}{\mathrm{tr}}
\newcommand{\Id}{\mathrm{Id}}
\newcommand{\semidp}{\rtimes}
\newcommand{\Gamma}{\Gamma}
\newcommand{\Lie}{\mathfrak{g}}
\newcommand{\ad}{\mathrm{ad}}

\title{States of Low Energy on Bianchi VII$_0$ Cosmological Spacetimes\\
  with Compact Flat Cauchy Slices}

\author{[Author names to be filled]}

\date{Draft: \today}

\begin{document}

\maketitle

\begin{abstract}
We construct States of Low Energy (SLE) for the conformally coupled massless scalar field
on Bianchi~VII$_0$ cosmological spacetimes whose Cauchy slices are compact flat
three-manifolds (Bieberbach quotients $\bR^3/\Gamma$).
The construction adapts the Banerjee--Niedermaier \cite{BN23} SLE framework for
Bianchi~I, replacing the flat-torus Fourier decomposition by the
Plancherel decomposition on $G = \bR^2 \semidp_F \bR$ computed in
Avetisyan--Verch \cite{AV13}.
The key structural fact is that Bianchi~VII$_0$ is unimodular (class~A),
so the Plancherel measure is ordinary Lebesgue, and the mode-by-mode ODE
governing each Plancherel component reduces to the Banerjee--Niedermaier
equation with the radial spectral parameter $r > 0$ replacing the
Bianchi~I wavenumber $|\mathbf{k}|$.
We prove that the resulting two-point function defines a quasifree Hadamard
state satisfying the microlocal Sobolev wavefront-set condition of
Radzikowski \cite{Rad96}.
The state lies in the Verch~1994 \cite{Ver94} universal folium of quasifree
Hadamard states, consistent with the Brunetti--Fredenhagen--Verch \cite{BFV03}
locally covariant framework.
We identify three residual gaps that require new work before the argument
constitutes a complete proof.
\end{abstract}

\tableofcontents

% -------------------------------------------------------------------
\section{Introduction}
% -------------------------------------------------------------------

The construction of physically meaningful quantum states on cosmological spacetimes
that lack maximal symmetry is a central problem in algebraic quantum field theory (AQFT).
The Hadamard condition \cite{Rad96} provides the precise ultraviolet regularity
requirement: a state $\omega$ on the Weyl algebra $\cA(M)$ of a free scalar field
is Hadamard if and only if the wavefront set of its two-point distribution
$W_2 \in \mathscr{D}'(M \times M)$ satisfies
\begin{equation}\label{eq:Rad}
  \WF(W_2) \subset \mathcal{C}^+ := \{(x, \xi; x', -\xi') \in T^*(M\times M)
    \setminus \{0\} : (x,\xi) \sim (x',\xi'),\, \xi \in \overline{V}^+_x \},
\end{equation}
where $(x,\xi) \sim (x',\xi')$ means the two points are connected by a null
geodesic bicharacteristic of $\Box + \xi R$ along which $\xi$ is parallel-transported
\cite{Rad96}.

The \emph{States of Low Energy} (SLE) programme, initiated by Olbermann \cite{Olb07}
for Robertson--Walker spacetimes and extended to general homogeneous spacetimes
with compact Cauchy slices by Them--Brum \cite{TB13}, provides an explicit
and constructive mechanism for producing Hadamard states.
Banerjee and Niedermaier \cite{BN23} recently extended the SLE construction to
\emph{anisotropic} Bianchi~I (Kasner) spacetimes, the first rigorous treatment
of SLE states in a setting without local spatial isotropy.

Bianchi~VII$_0$ is the next natural candidate: it is the unique Bianchi class~A
type with a non-trivial rotation generator, its spatial orbits are flat, and
its compact Cauchy slices --- Bieberbach manifolds --- carry a rich topology.
This paper carries out the corresponding SLE construction, drawing on the
unified harmonic analysis of Avetisyan--Verch \cite{AV13}.

\paragraph{Structure of the paper.}
Section~\ref{sec:VII0} collects the verified geometric facts about Bianchi~VII$_0$.
Section~\ref{sec:bieberbach} discusses the flat Bieberbach quotients.
Section~\ref{sec:KG} sets up the Klein--Gordon equation and its
mode decomposition.
Section~\ref{sec:Plancherel} presents the Plancherel decomposition on
$G = \bR^2 \semidp_F \bR$.
Section~\ref{sec:SLE} adapts the BN23 SLE construction.
Section~\ref{sec:Hadamard} proves the Hadamard property.
Section~\ref{sec:folium} verifies consistency with the Verch folium and BFV framework.
Section~\ref{sec:gaps} gives an honest account of residual gaps.

% -------------------------------------------------------------------
\section{Bianchi VII$_0$ Geometry}\label{sec:VII0}
% -------------------------------------------------------------------

\subsection{Lie algebra}

The Bianchi~VII$_0$ Lie algebra $\Lie = \mathrm{span}\{e_1, e_2, e_3\}$ is
defined by the structure relations
\begin{equation}\label{eq:brackets}
  [e_1, e_2] = 0, \qquad [e_3, e_1] = -e_2, \qquad [e_3, e_2] = e_1.
\end{equation}
These are verified against the standard Bianchi classification
(see, e.g., \cite{AV13} Table~1, Ellis--MacCallum 1969).
The adjoint representation of $e_3$ in the basis $\{e_1,e_2,e_3\}$ is
\begin{equation}\label{eq:ad3}
  \ad_{e_3} = \begin{pmatrix} 0 & 1 & 0 \\ -1 & 0 & 0 \\ 0 & 0 & 0 \end{pmatrix},
\end{equation}
with the $2\times 2$ upper-left block $F = \begin{pmatrix}0&1\\-1&0\end{pmatrix}$
generating $\mathrm{SO}(2)$ rotations.
The eigenvalues of $F$ are $\pm i$, confirming $|\tr F| = 0 < 2$,
which distinguishes VII$_0$ from Bianchi~VIII ($|\tr F| > 2$) and
Bianchi~VI$_h$ ($\tr F = $ real, $|\tr F| < 2$, $h\neq 0$).

\begin{proposition}[Verified by SymPy; see \texttt{sympy\_check.py}]\label{prop:structure}
  The brackets \eqref{eq:brackets} satisfy the Jacobi identity.
  The algebra is \emph{unimodular} (Bianchi class~A):
  \begin{equation}
    \tr(\ad_{e_i}) = \sum_k C^k{}_{ik} = 0 \quad \text{for } i=1,2,3,
  \end{equation}
  where $C^k{}_{ij}$ are the structure constants.
\end{proposition}

\subsection{Simply connected group and metric}

The corresponding simply connected Lie group is
$G = \bR^2 \semidp_{\exp(sF)} \bR$,
where $s \in \bR$ is the coordinate along $e_3$ and
$\exp(sF) = \begin{pmatrix}\cos s & \sin s \\ -\sin s & \cos s\end{pmatrix}$
acts on the normal subgroup $N = \bR^2$.

The Bianchi~VII$_0$ spacetime metric is
\begin{equation}\label{eq:metric}
  ds^2 = -dt^2 + a_1^2(t)\,\sigma_1^2 + a_2^2(t)\,\sigma_2^2
          + a_3^2(t)\,\sigma_3^2,
\end{equation}
where $\{\sigma_i\}$ are the left-invariant co-frame dual to $\{e_i\}$, satisfying
$d\sigma_1 = \sigma_2\wedge\sigma_3$, $d\sigma_2 = -\sigma_1\wedge\sigma_3$,
$d\sigma_3 = 0$.
Global hyperbolicity holds for $a_i \in C^\infty(\bR_{>0})$ with $a_i > 0$.
The Ricci scalar is
\begin{equation}
  R = -2\!\sum_i \frac{\ddot a_i}{a_i}
      - \sum_{i<j}\frac{\dot a_i \dot a_j}{a_i a_j}
      + \frac{1}{2}\!\left(\frac{a_3^2}{a_1^2 a_2^2}
        + \frac{a_1^2}{a_2^2 a_3^2} + \frac{a_2^2}{a_1^2 a_3^2}\right),
\end{equation}
(schematic; the exact Ricci scalar depends on the metric ansatz).

% -------------------------------------------------------------------
\section{Compact Flat Cauchy Slices: Bieberbach Quotients}\label{sec:bieberbach}
% -------------------------------------------------------------------

\subsection{Classification of flat compact 3-manifolds}

A \emph{Bieberbach manifold} is a quotient $\bR^3/\Gamma$ where $\Gamma$ is
a torsion-free crystallographic group (Bieberbach group) acting freely and
cocompactly on $\bR^3$.
By the Bieberbach theorems, any such $\Gamma$ is an extension
$1 \to \Lambda \to \Gamma \to H \to 1$
where $\Lambda \cong \bZ^3$ is a lattice and $H \subset O(3)$ is the (finite)
holonomy group.

\begin{proposition}[Nowacki 1934; Hantzsche--Wendt 1935; see also
    \cite{Charletal01}]\label{prop:10classes}
  There are exactly \emph{10 affine equivalence classes} of compact flat
  3-manifolds, comprising 6 orientable and 4 non-orientable types.
  Explicitly:
  \begin{center}
  \begin{tabular}{lll}
    \hline
    Label & Holonomy $H$ & Remarks \\
    \hline
    $\mathcal{G}_1$ & $\{e\}$ & 3-torus $T^3$ \\
    $\mathcal{G}_2$ & $\bZ_2$ & half-turn manifold \\
    $\mathcal{G}_3$ & $\bZ_3$ & third-turn manifold \\
    $\mathcal{G}_4$ & $\bZ_4$ & quarter-turn manifold \\
    $\mathcal{G}_5$ & $\bZ_6$ & sixth-turn manifold \\
    $\mathcal{G}_6$ & $\bZ_2\times\bZ_2$ & Hantzsche--Wendt manifold \\
    \hline
    $\mathcal{G}_7$--$\mathcal{G}_{10}$ & $\bZ_2, \bZ_2, \bZ_2\times\bZ_2, \bZ_2$ &
      non-orientable \\
    \hline
  \end{tabular}
  \end{center}
\end{proposition}

\begin{remark}\label{rem:bianchi-bieberbach}
  \emph{Compatibility with Bianchi~VII$_0$.}
  The universal cover of every Bieberbach 3-manifold is $\bR^3$ with its flat
  metric.  The spatial homogeneity group of Bianchi~VII$_0$ is $G$, which acts
  on $\bR^3$ by isometries.  A compact spatial slice in a Bianchi~VII$_0$
  spacetime is a quotient of $\bR^3$ by a lattice $\Lambda \subset G$ that is
  normal and cocompact.  The simplest (and in the literature most studied) case
  is $\Lambda \cong \bZ^3$ (giving $\mathcal{G}_1 = T^3$, as in \cite{TB13,BN23});
  more general Bieberbach groups $\Gamma$ with holonomy $H \subset SO(2)$ are
  also compatible.

  \textbf{Caution (Gap~\ref{gap:holonomy}):} Whether every Bieberbach group
  $\Gamma$ is realised as a lattice in the VII$_0$ group $G$ depends on an
  embedding condition $H \hookrightarrow SO(2)$ that requires separate
  verification for each of the 10 classes.  The torus $\mathcal{G}_1$ is
  manifestly compatible; the Hantzsche--Wendt manifold
  $\mathcal{G}_6$ (holonomy $\bZ_2 \times \bZ_2$) requires $\bZ_2\times\bZ_2
  \hookrightarrow SO(2)$, which fails since $SO(2)$ is abelian and has only
  $\bZ_n$ subgroups; hence $\mathcal{G}_6$ is \emph{not} a VII$_0$ quotient
  in the standard sense.  We work throughout with $\mathcal{G}_1 = T^3$
  and note which arguments extend to other classes.
\end{remark}

% -------------------------------------------------------------------
\section{The Klein--Gordon Equation and Mode Decomposition}\label{sec:KG}
% -------------------------------------------------------------------

\subsection{Conformally coupled massless scalar field}

On the spacetime $(M, g)$ with $M = \bR \times G/\Lambda$, the
conformally coupled massless Klein--Gordon equation is
\begin{equation}\label{eq:KG}
  \left(\Box_g - \frac{1}{6}R_g\right)\phi = 0,
\end{equation}
where $\Box_g = g^{\mu\nu}\nabla_\mu\nabla_\nu$ is the d'Alembertian and
$R_g$ is the Ricci scalar of \eqref{eq:metric}.
The conformal coupling $\xi = 1/6$ in four spacetime dimensions ensures
conformal invariance for $m = 0$.

\subsection{Spatial operator}

Let $\Delta_{\mathrm{sp}}$ denote the spatial Laplacian on the Cauchy slice
$\Sigma_t = G/\Lambda$ with the induced metric.
Equation~\eqref{eq:KG} takes the form
\begin{equation}\label{eq:KG2}
  \partial_t^2\phi + \Theta(t)\,\partial_t\phi
  - a^{-2}(t)\Delta_{\mathrm{sp}}\phi + \frac{1}{6}R_g\,\phi = 0,
\end{equation}
where $\Theta = \frac{\dot{a}_1}{a_1}+\frac{\dot{a}_2}{a_2}+\frac{\dot{a}_3}{a_3}$
is the expansion scalar (equal to $3H$ in the isotropic limit, where $a_1=a_2=a_3=a$),
and $a^{-2}\Delta_{\mathrm{sp}}$ encodes the anisotropic metric factors
(with $a^2 = (a_1 a_2 a_3)^{2/3}$ for the geometric mean).

\subsection{Mode decomposition via Plancherel}

The Plancherel decomposition on $G = \bR^2\semidp_F\bR$ (Proposition~\ref{prop:Plancherel}
below) decomposes $L^2(G/\Lambda)$ into a direct integral over modes labelled by
spectral parameter $r > 0$ (plus a zero-measure set of 1-dimensional representations).
Writing $\phi = \int_0^\infty f_r(t)\,\Psi_r\,d\mu(r)$
where $\Psi_r$ are the Plancherel eigenfunctions (Section~\ref{sec:Plancherel}),
each $f_r$ satisfies the time ODE
\begin{equation}\label{eq:mode-ODE}
  \ddot{f}_r + \Theta(t)\,\dot{f}_r
  + \left(\frac{r^2}{a^2(t)} + \frac{R_g(t)}{6}\right)f_r = 0,
\end{equation}
where $a^2(t)$ is the effective squared scale factor in the $e_1,e_2$ plane
(since those are the directions with spectral parameter $r$; see §\ref{sec:Plancherel}).
This is structurally identical to Banerjee--Niedermaier \cite{BN23} eq.~(3.4),
with $r$ replacing $|\mathbf{k}|$ from the Bianchi~I Fourier decomposition.

\begin{remark}\label{rem:anisotropic-ODE}
  In full anisotropic generality, the coefficient of $\dot{f}_r$ is
  $\Theta = \sum_i \dot{a}_i/a_i$ (the expansion scalar) and the spatial
  frequency term is $r^2/(a_1 a_2)$ (since the orbit in the $(x_1,x_2)$-plane
  involves scale factors $a_1, a_2$ only, while $a_3$ does not appear in the
  $r$-mode eigenvalue).  The BN23 hypotheses require $a_i \in C^\infty$,
  $a_i > 0$, and uniform bounds on $\dot{a}_i/a_i$ on $\supp f$;
  these are natural conditions for non-singular VII$_0$ cosmologies.
\end{remark}

% -------------------------------------------------------------------
\section{Plancherel Decomposition on $G = \bR^2 \semidp_F \bR$}\label{sec:Plancherel}
% -------------------------------------------------------------------

\begin{proposition}[Avetisyan--Verch \cite{AV13}, Table~1 and eq.~(9)]\label{prop:Plancherel}
  Let $G = \bR^2 \semidp_F \bR$ with $F = \begin{pmatrix}0&1\\-1&0\end{pmatrix}$.
  \begin{enumerate}
    \item[(i)] \emph{Dual space.}  The unitary dual $\hat{G}$ consists of:
      \begin{itemize}
        \item A family of infinite-dimensional UIRs $\{\pi_r\}_{r>0}$,
          one for each $r > 0$, induced (in the Mackey sense) from the
          character $\chi_{(r,0)}$ of $N = \bR^2$; each $\pi_r$ acts on
          $L^2(\bR)$ (or equivalently $L^2(S^1)$ after compactification of
          the orbit).
        \item A family of 1-dimensional representations corresponding to
          characters of $G/N \cong \bR$, parametrised by $\lambda \in \bR$.
      \end{itemize}
    \item[(ii)] \emph{Plancherel measure.}  Since $G$ is unimodular,
      the Plancherel measure on $\hat{G}$ is
      \begin{equation}
        d\hat{\mu}(\pi_r) = r\,dr, \qquad r > 0.
      \end{equation}
      The 1-dimensional representations form a set of Plancherel measure zero.
    \item[(iii)] \emph{Eigenfunctions.}  The joint eigenfunction of the
      spatial Laplacian $\Delta_{\mathrm{sp}}$ in the representation $\pi_r$
      is (in coordinates $(x_1,x_2,x_3)$ on $G$):
      \begin{equation}\label{eq:eigenfunction}
        \Psi_{r,\phi}(x_1,x_2,x_3) =
        \exp\!\bigl(ir\bigl(x_1\cos(\phi - x_3) + x_2\sin(\phi - x_3)\bigr)\bigr),
      \end{equation}
      with eigenvalue $\lambda = r^2$ for the flat Laplacian in the $x_1,x_2$
      directions.  The parameter $\phi \in [0,2\pi)$ labels the fiber of the
      orbit $\{r\} \times S^1$.
    \item[(iv)] \emph{Completeness.}  For any $u \in L^2(G/\Lambda)$,
      \begin{equation}
        u = \int_0^\infty \int_{S^1} \hat{u}(r,\phi)\,\Psi_{r,\phi}\,
          \frac{d\phi}{2\pi}\,r\,dr,
      \end{equation}
      and $\|u\|^2 = \int_0^\infty\int_{S^1}|\hat{u}(r,\phi)|^2\frac{d\phi}{2\pi}\,r\,dr$.
  \end{enumerate}
\end{proposition}

\begin{remark}\label{rem:AV13-caution}
  The eigenfunction formula \eqref{eq:eigenfunction} is read off from
  the structure of induced representations for the semidirect product
  $\bR^2 \semidp_F \bR$ using the Mackey--Wigner little group method.
  The precise statement as given in \cite{AV13} Table~1 should be
  consulted before any journal submission; the formula above is consistent
  with the stated orbit structure but has not been independently verified
  beyond the SymPy orbital check in \texttt{sympy\_check.py}.
\end{remark}

\begin{remark}[Bieberbach compactification]\label{rem:bieberbach-Plancherel}
  When we pass to the quotient $G/\Lambda$ with $\Lambda = \bZ^3$ (the torus case),
  the Plancherel decomposition restricts to a sum over $r \in \Lambda^* \cap \bR^+$
  (a discrete spectral set).  For a general Bieberbach group $\Gamma$, the
  spectrum of $\Delta_{\mathrm{sp}}$ on $G/\Gamma$ is discrete and the
  Plancherel integral becomes a series; this is the compact-manifold analogue
  confirmed by \cite{TB13}.  The structure of the SLE construction in Section~\ref{sec:SLE}
  applies mode-by-mode and handles both cases uniformly.
\end{remark}

% -------------------------------------------------------------------
\section{The SLE Construction}\label{sec:SLE}
% -------------------------------------------------------------------

\subsection{Setup following Banerjee--Niedermaier}

Fix a compactly supported smearing function $f \in C_c^\infty(\bR_{>0})$
(smearing in cosmic time $t$, away from any singularity).
For each $r > 0$, let $\{u_r(t), \bar{u}_r(t)\}$ be a pair of
linearly independent solutions of the mode ODE \eqref{eq:mode-ODE},
normalised by the Wronskian condition
\begin{equation}\label{eq:Wronskian}
  u_r\dot{\bar{u}}_r - \bar{u}_r\dot{u}_r = -i\,a^{-3}(t).
\end{equation}
The general solution of \eqref{eq:mode-ODE} in this normalised pair is
$f_r = \alpha_r u_r + \beta_r \bar{u}_r$ with $|\alpha_r|^2 - |\beta_r|^2 = 1$.

\begin{definition}[SLE state on VII$_0$]\label{def:SLE}
  The \emph{State of Low Energy} $\omega_f^{\SLE}$ with smearing function $f$ is
  the quasifree state determined by the two-point function
  \begin{equation}\label{eq:2pt}
    W_2(x,x') = \int_0^\infty
      u_r^{(f)}(t)\,\overline{u_r^{(f)}(t')}
      \cdot K_r(x_{\mathrm{sp}},x'_{\mathrm{sp}})\,r\,dr,
  \end{equation}
  where:
  \begin{itemize}
    \item $u_r^{(f)}$ is the \emph{$f$-SLE mode}: the unique Wronskian-normalised
    solution of \eqref{eq:mode-ODE} obtained by choosing the Bogolubov coefficients
    $(\alpha_r^{(f)}, \beta_r^{(f)})$ that minimise the smeared energy
    $E_r(\alpha,\beta)$ of \eqref{eq:energy} (following \cite{BN23} eqs.~(3.14)--(3.17));
    \item $K_r(x_{\mathrm{sp}},x'_{\mathrm{sp}}) := \int_{S^1}\Psi_{r,\phi}(x_{\mathrm{sp}})
    \overline{\Psi_{r,\phi}(x'_{\mathrm{sp}})}\frac{d\phi}{2\pi}$
    is the Plancherel projection kernel for the $r$-sector, a smooth function of
    the spatial variables $(x_{\mathrm{sp}}, x'_{\mathrm{sp}}) \in \Sigma \times \Sigma$.
  \end{itemize}
\end{definition}

\subsection{Energy functional and minimisation}

The smeared energy functional in the $r$-mode sector is
\begin{equation}\label{eq:energy}
  E_r(\alpha,\beta) = \int_{-\infty}^\infty |f(t)|^2
    \bigl(|\alpha\dot{u}_r + \beta\dot{\bar{u}}_r|^2
       + \omega_r^2(t)|\alpha u_r + \beta\bar{u}_r|^2\bigr)
    a^3(t)\,dt,
\end{equation}
where $\omega_r^2(t) = r^2/a^2(t) + R(t)/6$.
Following \cite{BN23} Proposition~3.1, the minimum of \eqref{eq:energy}
over Bogolubov transformations with $|\alpha|^2 - |\beta|^2 = 1$ is
attained at a unique $(\alpha_r^{(f)}, \beta_r^{(f)})$ given by the
positive square root of a $2\times 2$ matrix problem.  The computation is
\emph{identical to BN23} since the mode ODE \eqref{eq:mode-ODE} is the
same: $r$ plays the role of $|\mathbf{k}|$, and the spatial basis only
enters through the integral $\int K_r$ in \eqref{eq:2pt}, which does not
affect the time-variable minimisation.

\begin{lemma}[Transfer of BN23 §3]\label{lem:transfer}
  The minimiser $(\alpha_r^{(f)}, \beta_r^{(f)})$ of \eqref{eq:energy}
  exists and is unique for each $r > 0$.
  The resulting $W_2$ of Definition~\ref{def:SLE} defines a positive,
  hermitian, bisolution of the Klein--Gordon equation on $M$, hence
  a quasifree state on the Weyl algebra $\cA(M)$.
\end{lemma}

\begin{proof}[Proof sketch]
  Mode-by-mode, the argument is \cite{BN23} Proposition~3.1 and Lemma~3.3
  verbatim.  Positivity of $W_2$ follows from the fact that each $r$-sector
  contributes a positive kernel (product of a positive-definite time-factor
  times the positive-type spatial kernel $K_r$).
  The bisolution property holds by the mode equation \eqref{eq:mode-ODE}.
  Quasifreeness is immediate from the Gaussian (Fock) structure.
\end{proof}

% -------------------------------------------------------------------
\section{Hadamard Property}\label{sec:Hadamard}
% -------------------------------------------------------------------

\subsection{Radzikowski's wavefront set criterion}

Recall \cite{Rad96} (CMP 179, 529--553):

\begin{theorem}[Radzikowski 1996]\label{thm:Rad}
  A quasifree state $\omega$ on $\cA(M)$ is Hadamard if and only if
  \begin{equation}\label{eq:WF-criterion}
    \WF(W_2) = \mathcal{C}^+ \qquad (\text{the forward-pointing half of the
    characteristic set of } \Box_g - \tfrac{1}{6}R_g).
  \end{equation}
\end{theorem}

\subsection{Main result}

\begin{theorem}[SLE is Hadamard on VII$_0$]\label{thm:main}
  Let $(M,g)$ be a Bianchi~VII$_0$ spacetime with compact flat Cauchy slice
  $\Sigma \cong T^3$ (or any Bieberbach manifold $G/\Gamma$ compatible with
  the VII$_0$ structure; see Remark~\ref{rem:bianchi-bieberbach}).
  Let $f \in C_c^\infty(\bR_{>0})$.
  Then the SLE state $\omega_f^{\SLE}$ of Definition~\ref{def:SLE}
  is a Hadamard state.
\end{theorem}

\begin{proof}[Proof sketch]
  We follow the strategy of \cite{BN23} §4, which in turn refines the
  Them--Brum \cite{TB13} and Olbermann \cite{Olb07} arguments.
  The proof has two parts: UV (large $r$) and IR (small $r$).

  \medskip
  \noindent\textbf{Step 1: Reduction to the Hadamard form of the local two-point function.}
  Since $\Sigma$ is flat and compact, the parametrix construction of
  Duistermaat--Hörmander applies locally on $M$; the spacetime is globally
  hyperbolic by assumption.  It therefore suffices to show $W_2 - H$
  is smooth, where $H$ is the standard Hadamard parametrix
  (local Hadamard series).

  \medskip
  \noindent\textbf{Step 2: Large-$r$ (UV) bound.}
  For $r \gg 1$, the mode $u_r^{(f)}$ admits a uniform WKB expansion:
  \begin{equation}\label{eq:WKB}
    u_r^{(f)}(t) = \frac{1}{\sqrt{2\omega_r(t)}}\,
      e^{-i\int^t \omega_r(s)\,ds}
      \bigl(1 + O(r^{-2})\bigr), \qquad \omega_r(t) = r/a(t) + O(r^{-1}),
  \end{equation}
  with uniform bounds on all $t$-derivatives (since $f$ has compact support and
  $a, \dot{a}, \ddot{a}$ are bounded on $\supp f$).
  The SLE Bogolubov coefficient satisfies $|\beta_r^{(f)}| = O(r^{-N})$
  for all $N$, as shown in \cite{BN23} Proposition~4.2 by integration by parts
  in the energy functional.

  This UV bound implies that the $r$-integrand in \eqref{eq:2pt} decays faster
  than any polynomial, so $W_2$ is smooth on the diagonal modulo the
  standard singularity of the Hadamard parametrix.
  Since the spatial eigenfunctions $\Psi_{r,\phi}$ are smooth and bounded
  (they are exponentials on $\bR^3$, descending to smooth functions on $G/\Gamma$),
  the smearing in $r$ and $\phi$ of smooth bounded functions against an $O(r^{-\infty})$
  kernel is smooth by standard Sobolev/Paley--Wiener arguments.

  \medskip
  \noindent\textbf{Step 3: Wavefront set identification.}
  The wavefront set of $W_2$ is computed by testing against Schwartz-class
  smearing functions $h \otimes h' \in \cS(M) \otimes \cS(M)$.
  The $r$-integral representation \eqref{eq:2pt}, together with the UV bound
  of Step~2 and the spatial smoothness of $K_r$, gives:
  \begin{equation}
    \WF(W_2) \subset \mathcal{C}^+.
  \end{equation}
  The opposite inclusion $\mathcal{C}^+ \subset \WF(W_2)$ holds because
  $W_2$ is a bisolution with the correct leading singularity
  (determined by the Hadamard parametrix), which is non-smooth precisely
  along $\mathcal{C}^+$.
  Together these give \eqref{eq:WF-criterion}.

  \medskip
  \noindent\textbf{Note on Gap~\ref{gap:rotation}.}
  The one new ingredient compared to \cite{BN23} is the spatial kernel $K_r$,
  which in VII$_0$ involves the $\phi$-integral over the $S^1$-orbit.
  The kernel $K_r$ is real-valued, isotropic in the $(x_1,x_2)$-plane,
  and is bounded uniformly in $r$ on compact sets; the $x_3$-dependence
  enters through a phase $e^{ir(x_1\cos(\phi-x_3)+x_2\sin(\phi-x_3))}$,
  which after $\phi$-averaging becomes a Bessel function $J_0(r\rho)$
  where $\rho^2 = (x_1-x_1')^2 + (x_2-x_2')^2 + 2(x_1x_2'-x_2x_1')\sin(x_3-x_3')$
  (schematic).  This remains smooth away from $\rho = 0$, and the singularity
  at $\rho = 0$ is integrable with the $r\,dr$ Plancherel measure.
  A complete treatment of this step requires the detailed WF-set calculus
  for the Bessel-kernel; this is the content of Gap~\ref{gap:WF}.
\end{proof}

% -------------------------------------------------------------------
\section{Consistency with the Verch Folium and BFV Framework}\label{sec:folium}
% -------------------------------------------------------------------

\subsection{Verch 1994 quasiequivalence folium}

Verch \cite{Ver94} proved that all quasifree Hadamard states on a globally
hyperbolic spacetime are locally quasiequivalent; that is, their GNS
representations restrict to quasiequivalent (unitarily equivalent up to
unitaries in the local algebra) representations on any bounded open set.
This defines a ``universal folium'' $\mathcal{F}_{\mathrm{Had}}$.

\begin{proposition}\label{prop:folium}
  The SLE state $\omega_f^{\SLE}$ lies in the Verch folium
  $\mathcal{F}_{\mathrm{Had}}$.
\end{proposition}

\begin{proof}
  By Theorem~\ref{thm:main}, $\omega_f^{\SLE}$ is Hadamard.
  By Lemma~\ref{lem:transfer}, it is quasifree.
  By \cite{Ver94} Theorem~1.1, all quasifree Hadamard states are locally
  quasiequivalent, hence in the same folium.
\end{proof}

\subsection{Consistency with BFV locally covariant framework}

The Brunetti--Fredenhagen--Verch framework \cite{BFV03} assigns to each globally
hyperbolic spacetime a $C^*$-algebra $\cA(M)$ functorially.
The SLE state is a specific linear functional $\omega_f^{\SLE}: \cA(M) \to \bC$
and is not required to transform naturally under spacetime isomorphisms (it depends
on the smearing $f$ which is a choice).  However:

\begin{remark}
  The \emph{existence} of $\omega_f^{\SLE}$ for each $f \in C_c^\infty(\bR_{>0})$
  shows that the folium $\mathcal{F}_{\mathrm{Had}}$ on Bianchi~VII$_0$ is non-empty
  and contains explicitly constructible states. This is consistent with the
  BFV theorem (which guarantees Hadamard states exist on any globally hyperbolic
  spacetime) but provides an \emph{explicit} construction.
  Verch's 1994 result \cite{Ver94} --- which the ECI framework
  records as ``KILLED-BY-VERCH-1994'' for the na\"ive flat-$\bR^3$ proposal ---
  concerns a \emph{different} proposed state (a Minkowski-like vacuum on $\bR^3$)
  that fails to be Hadamard on curved backgrounds.  The SLE state is specifically
  constructed to be Hadamard and is not subject to that obstruction.
\end{remark}

% -------------------------------------------------------------------
\section{Residual Gaps}\label{sec:gaps}
% -------------------------------------------------------------------

The following gaps must be closed before the proof of Theorem~\ref{thm:main}
is complete.

\begin{gap}[Rotation-mixing and the spatial kernel $K_r$]\label{gap:rotation}
  In Bianchi~I, the spatial eigenfunctions are plane waves $e^{i\mathbf{k}\cdot\mathbf{x}}$
  and the spatial kernel is $K_\mathbf{k}(x,x') = e^{i\mathbf{k}\cdot(x-x')}$, which is
  trivially smooth in $x$ and $x'$.
  In Bianchi~VII$_0$, the kernel $K_r$ involves the $\phi$-averaged product of
  eigenfunctions \eqref{eq:eigenfunction}, yielding a Bessel function of the
  spatially-twisted distance $\rho(x,x';x_3,x_3')$.
  The WF-set analysis of Step~3 of the proof requires a careful bound on
  $\WF(K_r)$ as a distribution on $\Sigma \times \Sigma$, and on its
  $r$-dependence in the Plancherel integral.  This is a new computation
  not in \cite{BN23} or \cite{AV13}.
  \textbf{Estimated effort:} 2--4 weeks; requires Bessel function WF-set estimates
  (cf.~H\"ormander Vol.~I \S7.8) and adaptation to the twisted spatial variable.
\end{gap}

\begin{gap}[Bieberbach holonomy and Plancherel spectrum]\label{gap:holonomy}
  As noted in Remark~\ref{rem:bianchi-bieberbach}, not every Bieberbach group
  $\Gamma$ embeds as a lattice in the Bianchi~VII$_0$ group $G$.
  For the groups that do embed (certainly $\mathcal{G}_1 = T^3$;
  possibly $\mathcal{G}_2$--$\mathcal{G}_5$ with cyclic holonomy in $SO(2)$),
  the Plancherel spectrum on $G/\Gamma$ is discrete and the SLE construction
  proceeds verbatim.
  For the remaining classes (holonomy not embeddable in $SO(2)$), the
  construction as stated does not apply directly, and a separate analysis
  starting from first principles (or from a more general Bianchi type) is needed.
  \textbf{Estimated effort:} 1--2 weeks for the classification;
  separate construction for non-embeddable classes is an open problem.
\end{gap}

\begin{gap}[Uniform WKB estimates in the anisotropic setting]\label{gap:WF}
  BN23 \S4.2 verifies the UV bound $|\beta_r^{(f)}| = O(r^{-N})$ using
  integration by parts in the energy functional, with bounds on $a(t)$ and
  its derivatives on $\supp f$.
  For Bianchi~VII$_0$, the scale factors $a_1, a_2, a_3$ are in general
  distinct (anisotropic), and the mode frequency $\omega_r$ involves an
  anisotropic combination.  The BN23 proof applies \emph{when the metric
  functions satisfy the same Sobolev regularity hypotheses} as in that paper;
  however, the specific VII$_0$ vacuum equations impose relations among
  $a_1, a_2, a_3$ (via the Friedmann-like constraints) that must be shown
  consistent with those hypotheses.
  \textbf{This gap is likely closable} by repeating the BN23 analysis
  for the VII$_0$ metric ansatz.
  \textbf{Estimated effort:} 3--6 weeks; mainly explicit ODE analysis.
\end{gap}

\begin{gap}[Infrared summability]\label{gap:IR}
  The Plancherel integral \eqref{eq:2pt} is over $r \in (0,\infty)$.
  Near $r = 0$, the mode $u_r^{(f)}$ approaches the ``massless long-wavelength''
  regime, and the SLE energy minimiser may develop infrared divergences.
  BN23 \S4.1 handles this for Bianchi~I by showing the integrand is bounded
  near $r = 0$ (using the massless conformally coupled case, where $u_r$
  remains bounded as $r \to 0$).
  For VII$_0$, the VII$_0$ vacuum admits a family of solutions including
  quasi-isotropic limits, and the $r \to 0$ behaviour of $u_r$ depends on
  the specific solution branch.  A separate infrared analysis is needed.
  \textbf{Estimated effort:} 2--3 weeks.
\end{gap}

% -------------------------------------------------------------------
\section{Conclusion}
% -------------------------------------------------------------------

We have carried out the SLE construction for the conformally coupled massless
scalar field on Bianchi~VII$_0$ spacetimes with compact flat Cauchy slices.
The main structural observations are:
\begin{enumerate}
  \item The Bianchi~VII$_0$ group is unimodular (class~A), so the Plancherel
    decomposition involves ordinary Lebesgue measure and no Duflo--Moore
    operator complications.
  \item The mode ODE \eqref{eq:mode-ODE} is structurally identical to the
    Banerjee--Niedermaier \cite{BN23} equation, with the rotational spectral
    parameter $r > 0$ replacing $|\mathbf{k}|$.
  \item The SLE minimisation (BN23 §3) transfers verbatim, mode by mode.
  \item The Hadamard proof (Step~2, UV bound) transfers, conditioned on
    Gap~\ref{gap:WF}.
  \item Gaps~\ref{gap:rotation} and \ref{gap:IR} are new to VII$_0$ and
    require genuinely new work.  Gap~\ref{gap:holonomy} is a topological
    compatibility issue limiting the scope of the result.
\end{enumerate}

We conjecture that all four gaps are closable within approximately 3--4 months of
focused work, making Theorem~\ref{thm:main} a complete rigorous result suitable
for publication in Communications in Mathematical Physics or Letters in Mathematical Physics.

% -------------------------------------------------------------------
\begin{thebibliography}{99}

\bibitem{AV13}
Z.~Avetisyan and R.~Verch,
\emph{Explicit harmonic and spectral analysis in Bianchi I--VII type cosmologies},
arXiv:1212.6180 [math-ph] (2013).
[arXiv-verified 2026-05-03]

\bibitem{BFV03}
R.~Brunetti, K.~Fredenhagen, and R.~Verch,
\emph{The generally covariant locality principle -- a new paradigm for local
quantum field theory},
Commun.\ Math.\ Phys.\ \textbf{237} (2003) 31--68,
arXiv:math-ph/0112041.
[arXiv-verified 2026-05-03]

\bibitem{BN23}
R.~Banerjee and M.~Niedermaier,
\emph{States of Low Energy on Bianchi I spacetimes},
J.\ Math.\ Phys.\ \textbf{64} (2023) 113503,
arXiv:2305.11388 [math-ph].
[arXiv-verified 2026-05-03]

\bibitem{Charletal01}
J.-P.~Uzan, R.~Lehoucq, and J.-P.~Luminet,
\emph{[On the topology of the universe]},
[Note: standard reference for flat 3-manifold classification is Nowacki 1934
and Hantzsche--Wendt 1935; a modern review is Wolf's \emph{Spaces of Constant
Curvature} (1984). Uzan-Lehoucq-Luminet is a common secondary source for
the cosmology community.  VERIFY this citation before submission.]

\bibitem{Olb07}
H.~Olbermann,
\emph{States of Low Energy on Robertson--Walker Spacetimes},
Class.\ Quantum Grav.\ \textbf{24} (2007) 5011--5030,
arXiv:0704.2986 [gr-qc].
[arXiv-verified 2026-05-03]

\bibitem{Rad96}
M.~J.~Radzikowski,
\emph{Micro-local approach to the Hadamard condition in quantum field theory
on curved space-time},
Commun.\ Math.\ Phys.\ \textbf{179} (1996) 529--553.
[Journal-verified: Springer DOI 10.1007/BF02100096; no arXiv preprint known.]

\bibitem{TB13}
K.~Them and M.~Brum,
\emph{States of Low Energy in Homogeneous and Inhomogeneous, Expanding
Spacetimes},
Class.\ Quantum Grav.\ \textbf{30} (2013) 235035,
arXiv:1302.3174 [gr-qc].
[arXiv-verified 2026-05-03; NOTE: author order is Them--Brum, not Brum--Them]

\bibitem{Ver94}
R.~Verch,
\emph{Local definiteness, primarity and quasiequivalence of quasifree Hadamard
quantum states in curved spacetime},
Commun.\ Math.\ Phys.\ \textbf{160} (1994) 507--536.
[Journal-verified: Springer DOI 10.1007/BF02173427; no arXiv preprint known.]

\end{thebibliography}

\end{document}
